\documentclass[12pt, a4paper]{report}

% ── packages ──────────────────────────────────────────────────────────────────
\usepackage[a4paper, margin=2.5cm]{geometry}
\usepackage[T1]{fontenc}
\usepackage[utf8]{inputenc}


\usepackage{xcolor}
\usepackage{listings}
\usepackage{titlesec}
\usepackage{titletoc}
\usepackage{fancyhdr}
\usepackage{booktabs}
\usepackage{amsmath}
\usepackage{colortbl}
\usepackage{array}
\usepackage{tabularx}
\usepackage{enumitem}
\usepackage{mdframed}
\usepackage{parskip}
\usepackage{graphicx}
\usepackage{hyperref}
\usepackage{tocloft}
\usepackage{setspace}

% ── colours ───────────────────────────────────────────────────────────────────
\definecolor{accent}{HTML}{4F46E5}
\definecolor{accentlight}{HTML}{EEF2FF}
\definecolor{darkgrey}{HTML}{2D2D2D}
\definecolor{midgrey}{HTML}{555555}
\definecolor{lightgrey}{HTML}{F2F2F2}
\definecolor{codebg}{HTML}{F8F8F8}
\definecolor{border}{HTML}{D1D5DB}
\definecolor{noteblue}{HTML}{EFF6FF}
\definecolor{noteborder}{HTML}{3B82F6}
\definecolor{stringcol}{HTML}{16A34A}
\definecolor{commentcol}{HTML}{6B7280}
\definecolor{keywordcol}{HTML}{DC2626}
\definecolor{numbercol}{HTML}{D97706}

% ── hyperref setup ────────────────────────────────────────────────────────────
\hypersetup{
  colorlinks=true,
  linkcolor=accent,
  urlcolor=accent,
  citecolor=accent,
  pdftitle={Custom Memory Allocator},
  pdfauthor={Samir Mansour},
  pdfsubject={Systems Programming Portfolio Report},
}

% ── code listings ─────────────────────────────────────────────────────────────
\lstdefinestyle{cstyle}{
  language=C,
  backgroundcolor=\color{codebg},
  basicstyle=\ttfamily\small,
  keywordstyle=\color{keywordcol}\bfseries,
  stringstyle=\color{stringcol},
  commentstyle=\color{commentcol}\itshape,
  numberstyle=\tiny\color{numbercol},
  numbers=left,
  numbersep=8pt,
  stepnumber=1,
  showstringspaces=false,
  breaklines=true,
  breakatwhitespace=false,
  frame=leftline,
  framesep=4pt,
  rulecolor=\color{accent},
  framerule=3pt,
  xleftmargin=14pt,
  xrightmargin=4pt,
  tabsize=4,
  captionpos=b,
  aboveskip=10pt,
  belowskip=6pt,
}

\lstdefinestyle{shellstyle}{
  backgroundcolor=\color{darkgrey!95!black},
  basicstyle=\ttfamily\small\color{white},
  showstringspaces=false,
  breaklines=true,
  frame=single,
  framesep=6pt,
  rulecolor=\color{accent},
  framerule=1pt,
  xleftmargin=4pt,
  xrightmargin=4pt,
  aboveskip=8pt,
  belowskip=6pt,
}

\lstset{style=cstyle}

% ── section styling ───────────────────────────────────────────────────────────
\titleformat{\chapter}[block]
  {\huge\bfseries\color{accent}}
  {\thechapter.}{0.6em}{}
  [\vspace{2pt}\color{accent}\rule{\linewidth}{1.5pt}]

\titleformat{\section}[block]
  {\large\bfseries\color{darkgrey}}
  {\thesection}{0.5em}{}

\titleformat{\subsection}[block]
  {\normalsize\bfseries\color{midgrey}}
  {\thesubsection}{0.5em}{}

\titlespacing*{\chapter}{0pt}{0pt}{16pt}
\titlespacing*{\section}{0pt}{14pt}{6pt}
\titlespacing*{\subsection}{0pt}{10pt}{4pt}

% ── TOC styling ───────────────────────────────────────────────────────────────
\renewcommand{\cftchapfont}{\bfseries\color{accent}}
\renewcommand{\cftchappagefont}{\bfseries\color{accent}}
\renewcommand{\cftsecfont}{\color{darkgrey}}
\renewcommand{\cftsecpagefont}{\color{darkgrey}}
\renewcommand{\cftsubsecfont}{\small\color{midgrey}}
\renewcommand{\cftsubsecpagefont}{\small\color{midgrey}}
\setlength{\cftbeforechapskip}{8pt}
\setlength{\cftbeforesecskip}{2pt}

% ── header and footer ─────────────────────────────────────────────────────────
\pagestyle{fancy}
\fancyhf{}
\fancyhead[L]{\small\color{midgrey}Custom Memory Allocator}
\fancyhead[R]{\small\color{midgrey}Samir Mansour -- github.com/Samir204}
\fancyfoot[C]{\small\color{midgrey}\thepage}
\renewcommand{\headrulewidth}{0.4pt}
\renewcommand{\headrule}{\hbox to\headwidth{\color{border}\leaders\hrule height \headrulewidth\hfill}}

% ── note box ─────────────────────────────────────────────────────────────────
\newmdenv[
  backgroundcolor=noteblue,
  linecolor=noteborder,
  linewidth=3pt,
  leftline=true,
  rightline=false,
  topline=false,
  bottomline=false,
  innerleftmargin=10pt,
  innerrightmargin=10pt,
  innertopmargin=8pt,
  innerbottommargin=8pt,
  skipabove=8pt,
  skipbelow=8pt,
]{notebox}

% ── list styling ──────────────────────────────────────────────────────────────
\setlist[itemize]{
  leftmargin=1.5em,
  itemsep=3pt,
  parsep=0pt,
  topsep=4pt,
}
\setlist[itemize,1]{label=\textcolor{accent}{\textbullet}}
\setlist[itemize,2]{label=\textcolor{midgrey}{\textendash}}

% ── table styling ─────────────────────────────────────────────────────────────
\newcolumntype{L}[1]{>{\raggedright\arraybackslash}p{#1}}
\newcolumntype{C}[1]{>{\centering\arraybackslash}p{#1}}

\arrayrulecolor{border}
\renewcommand{\arraystretch}{1.35}

% ── misc ──────────────────────────────────────────────────────────────────────
\setlength{\parindent}{0pt}
\setlength{\parskip}{8pt}
\onehalfspacing

% ─────────────────────────────────────────────────────────────────────────────
\begin{document}

% ── COVER PAGE ────────────────────────────────────────────────────────────────
\begin{titlepage}
  \pagecolor{accent}
  \color{white}
  \vspace*{3cm}

  {\fontsize{36}{44}\selectfont\bfseries Custom Memory\\[6pt] Allocator\par}

  \vspace{1cm}
  \color{white!80!accent}
  {\Large malloc / free / realloc from scratch in C\par}

  \vspace{0.6cm}
  \color{white}
  \rule{\linewidth}{0.5pt}
  \vspace{0.6cm}

  \begin{tabular}{ll}
    \textbf{Author}   & Samir Mansour \\[4pt]
    \textbf{GitHub}   & github.com/Samir204 \\[4pt]
    \textbf{Project}  & Summer Portfolio -- Systems Programming \\[4pt]
    \textbf{Date}     & July 2026 \\[4pt]
    \textbf{Language} & C / C++ \\
  \end{tabular}

  \vfill

  \color{white!60!accent}
  \rule{\linewidth}{0.5pt}
  \vspace{0.3cm}

  {\small Memory Management \quad|\quad Heap Allocator \quad|\quad Benchmarking \quad|\quad Thread Safety\par}
\end{titlepage}

\nopagecolor

% ── TABLE OF CONTENTS ─────────────────────────────────────────────────────────
\tableofcontents
\thispagestyle{empty}
\clearpage

% ─────────────────────────────────────────────────────────────────────────────
\chapter{Introduction}
% ─────────────────────────────────────────────────────────────────────────────

When you write \texttt{malloc(32)} in a C program, the standard library does a
surprising amount of work behind the scenes. It finds a free chunk of heap
memory, tracks it, hands you a pointer, and later takes it back when you call
\texttt{free()}. This project builds all of that machinery from scratch in C,
without using any of the standard library allocator functions.

The goal is not to beat glibc in a benchmark. The goal is to understand exactly
what glibc is doing, why it makes the trade-offs it does, and to be able to talk
about memory management at a level most programmers never reach. That depth is
what separates someone who \textit{knows C} from someone who actually understands
how computers work underneath.

This report documents every part of the implementation: the data structures, all
three allocation strategies (first-fit, best-fit, and buddy system), thread
safety, corruption detection, and the benchmarking results that show the
trade-offs between each approach.

\section{What This Project Implements}

\begin{itemize}
  \item A complete heap allocator using \texttt{sbrk} and \texttt{mmap} system calls
  \item Three allocation strategies: first-fit, best-fit, and buddy system
  \item Block splitting to avoid wasting memory inside large free blocks
  \item Block coalescing to merge adjacent free blocks and reduce fragmentation
  \item Thread safety using a POSIX mutex
  \item Double-free and invalid pointer detection using magic number headers
  \item In-place \texttt{realloc} growth when the next block is free and adjacent
  \item A benchmarking harness comparing throughput and fragmentation ratio across all three strategies
\end{itemize}

\section{Repository Structure}

\begin{lstlisting}[style=shellstyle, numbers=none]
allocator/
  allocator.h       public API and struct definitions
  allocator.c       main implementation (first-fit, best-fit)
  buddy.c           buddy system allocator
  test.c            unit test suite
  benchmark.c       benchmarking harness
  results.csv       benchmark output data
  Makefile          build configuration
\end{lstlisting}

% ─────────────────────────────────────────────────────────────────────────────
\chapter{Background and Motivation}
% ─────────────────────────────────────────────────────────────────────────────

Memory management is one of the most fundamental problems in systems
programming. Every non-trivial program needs to allocate memory at runtime, and
how well it does so directly affects performance, reliability, and security.

\section{How the Heap Works}

When a process starts, the operating system gives it a virtual address space.
The \textbf{heap} is a region of this space that can grow upward as the program
needs more memory. The boundary between the in-use heap and the unmapped region
above it is called the \textbf{program break}.

There are two main ways a process can ask the OS for more heap memory:

\begin{itemize}
  \item \texttt{sbrk(n)}: moves the program break forward by \texttt{n} bytes,
    returning the old break address. Simple and sequential, but limited to one
    contiguous region that can only grow at the end.
  \item \texttt{mmap}: maps a new anonymous region of memory anywhere in the
    address space. More flexible, works in page-sized chunks (4096 bytes each),
    and individual mappings can be returned to the OS independently.
\end{itemize}

This allocator uses \texttt{sbrk} for the main heap (first-fit and best-fit
strategies) and a fixed pool for the buddy system. This mirrors how real
allocators like glibc work: \texttt{sbrk} for small allocations, \texttt{mmap}
for large ones.

\section{Why Not Just Call sbrk Every Time}

You cannot call \texttt{sbrk} for every individual \texttt{malloc}. A program
that does a million small allocations would make a million system calls, which
would be catastrophically slow. The allocator's job is to call \texttt{sbrk}
infrequently to grab large chunks of memory, then manage the subdivision and
reuse of those chunks in userspace. That management is what makes a memory
allocator non-trivial to write correctly.

\section{Fragmentation}

Fragmentation is the core problem that makes allocator design hard. It comes in
two forms:

\begin{itemize}
  \item \textbf{External fragmentation}: free memory exists but is split into
    pieces too small to satisfy a request, even though the total free space would
    be enough in aggregate.
  \item \textbf{Internal fragmentation}: a block is larger than the requested
    size, so some bytes inside it are wasted and unavailable to anyone else.
\end{itemize}

Every allocation strategy makes a different trade-off between these two forms of
fragmentation, and between fragmentation and throughput. Measuring that
trade-off is the main deliverable of Week 4.

% ─────────────────────────────────────────────────────────────────────────────
\chapter{Data Structures}
% ─────────────────────────────────────────────────────────────────────────────

\section{The Block Header}

Every chunk of memory managed by the allocator has a small header placed
immediately before the user data. The allocator uses this header to track each
block without any external bookkeeping:

\begin{lstlisting}
typedef struct Block {
    size_t   size;      // usable bytes after this header
    int      is_free;   // 1 = available, 0 = in use
    uint32_t magic;     // corruption and double-free detection
    struct Block *next; // next block in linked list
    struct Block *prev; // previous block
} Block;
\end{lstlisting}

On a 64-bit system this struct is 32 bytes. The layout in memory looks like:

\begin{lstlisting}[style=shellstyle, numbers=none]
[ header 32B ][ user data ][ header 32B ][ user data ] ...
\end{lstlisting}

Given a user pointer \texttt{ptr}, recovering the header is always:

\begin{lstlisting}
Block *header = (Block*)ptr - 1;
\end{lstlisting}

And returning the user pointer from a header is always:

\begin{lstlisting}
void *user_ptr = (void*)(header + 1);
\end{lstlisting}

This single pattern is used everywhere in the allocator. Getting it wrong by
even one byte silently corrupts memory in ways that are nearly impossible to
debug after the fact.

\section{Magic Numbers}

The \texttt{magic} field stores a constant that identifies a valid block header.
Two values are used:

\begin{lstlisting}
#define MAGIC_ALLOC 0xDEADBEEF   // block is currently allocated
#define MAGIC_FREE  0xDEADDEAD   // block has been freed
\end{lstlisting}

Every new block gets \texttt{MAGIC\_ALLOC} stamped on creation, and again
whenever it is reused by \texttt{my\_malloc}. When \texttt{my\_free} is called,
it checks the magic field before touching anything else. If it sees
\texttt{MAGIC\_FREE}, the block was already freed -- the call is rejected with
an error message. If it sees anything other than \texttt{MAGIC\_ALLOC}, the
pointer is invalid and also rejected. Neither case crashes the program, which
makes debugging far easier.

\begin{notebox}
A critical bug discovered during development: when \texttt{my\_malloc} reuses a
freed block, the magic field must be reset to \texttt{MAGIC\_ALLOC}. Missing
this caused every subsequent \texttt{free} on that block to be reported as a
double-free, which produced hundreds of false error messages during the thread
safety test.
\end{notebox}

\section{The Heap Linked List}

All blocks on the heap are linked in a doubly linked list starting at
\texttt{heap\_start}. The list is ordered by address, since blocks are always
appended to the end when \texttt{extend\_heap} is called. Walking this list is
how \texttt{find\_free} and \texttt{find\_best\_fit} search for available
blocks.

The linked list approach is simple and correct but O(n) for allocation.
Production allocators use segregated free lists (one list per size class) to
achieve O(1) or O(log n) allocation. This is the main area where glibc
outperforms this implementation on large heaps.

% ─────────────────────────────────────────────────────────────────────────────
\chapter{Allocation Strategies}
% ─────────────────────────────────────────────────────────────────────────────

\section{First-Fit}

First-fit scans the block list from the start and returns the first free block
large enough to satisfy the request:

\begin{lstlisting}
static Block *find_free(size_t size) {
    Block *cur = heap_start;
    while (cur) {
        if (cur->is_free && cur->size >= size)
            return cur;
        cur = cur->next;
    }
    return NULL;
}
\end{lstlisting}

First-fit is fast because it stops at the first match. The downside is that it
tends to fill the beginning of the heap with small leftover fragments, causing
subsequent searches to slow down over time as they must skip over many small
used blocks before finding a suitable free one.

\section{Best-Fit}

Best-fit scans the entire list and returns the smallest block that is still
large enough for the request:

\begin{lstlisting}
static Block *find_best_fit(size_t size) {
    Block *cur = heap_start;
    Block *best = NULL;
    while (cur) {
        if (cur->is_free && cur->size >= size) {
            if (best == NULL || cur->size < best->size)
                best = cur;
        }
        cur = cur->next;
    }
    return best;
}
\end{lstlisting}

Best-fit reduces internal fragmentation because it finds the tightest match,
leaving larger free blocks available for future requests. The cost is that it
always does a full list scan, making it consistently slower than first-fit.
Interestingly, best-fit can actually produce \textit{more} external fragmentation
than first-fit on certain workloads, because choosing the tightest match tends
to leave many tiny unusable remainders.

\section{Block Splitting}

Both first-fit and best-fit use block splitting. When a free block is
significantly larger than needed, it is split into two: an exact-fit allocated
block and a smaller free remainder. Without splitting, a 1024-byte free block
used to satisfy an 8-byte request would waste 1016 bytes permanently.

\begin{lstlisting}
#define MIN_SPLIT (HEADER_SIZE + 8)

static void split_block(Block *block, size_t size) {
    if (block->size < size + MIN_SPLIT) return;

    Block *new_block = (Block*)((char*)(block + 1) + size);
    new_block->size    = block->size - size - HEADER_SIZE;
    new_block->is_free = 1;
    new_block->magic   = MAGIC_ALLOC;
    new_block->next    = block->next;
    new_block->prev    = block;

    if (block->next) block->next->prev = new_block;
    block->next = new_block;
    block->size = size;
}
\end{lstlisting}

\texttt{MIN\_SPLIT} is set to \texttt{HEADER\_SIZE + 8}. A block is only split
if the remainder would be at least large enough to hold a header plus 8 bytes of
usable data. Creating a block smaller than this would just add overhead without
any practical benefit.

\section{Coalescing}

When a block is freed, if its neighbours in memory are also free, they are
merged into one larger block. This is \textbf{coalescing}, and it is essential
for preventing the heap from filling up with unusable tiny fragments over time:

\begin{lstlisting}
void coalesce(Block *block) {
    // merge with next block if it is free
    if (block->next && block->next->is_free) {
        block->size += HEADER_SIZE + block->next->size;
        block->next  = block->next->next;
        if (block->next) block->next->prev = block;
    }
    // merge with previous block if it is free
    if (block->prev && block->prev->is_free) {
        block->prev->size += HEADER_SIZE + block->size;
        block->prev->next  = block->next;
        if (block->next) block->next->prev = block->prev;
    }
}
\end{lstlisting}

Coalescing is called at the end of every \texttt{my\_free}. The key constraint
is that it must only do pointer surgery on existing blocks. Calling
\texttt{my\_malloc} or \texttt{my\_free} inside coalesce would cause infinite
recursion.

\section{The Buddy System}

The buddy system is a fundamentally different approach. Instead of blocks being
arbitrary sizes, every block is always a power-of-two size. When memory is
needed, a large block is halved repeatedly until the smallest power-of-two that
fits the request is reached. When freed, a block checks whether its partner (the
\textbf{buddy}) is also free. If so, they merge back together. Then that merged
block checks its buddy, and so on up the chain.

The key insight is how to find a block's buddy. Given a block at address
\texttt{A} with size \texttt{S} and pool base address \texttt{base}:

\begin{lstlisting}[style=shellstyle, numbers=none]
buddy_address = base + ((A - base) XOR S)
\end{lstlisting}

XOR with the block size flips exactly the one bit that separates the block from
its buddy. This makes buddy-finding O(1) with no searching required:

\begin{lstlisting}
static BuddyBlock *get_buddy(BuddyBlock *block) {
    size_t offset      = (char*)block - buddy_pool;
    size_t buddy_offset = offset ^ block->size;
    if (buddy_offset >= BUDDY_POOL_SIZE) return NULL;
    return (BuddyBlock*)(buddy_pool + buddy_offset);
}
\end{lstlisting}

All allocation and deallocation operations in the buddy system are O(log n)
where n is the number of size levels. The trade-off is internal fragmentation: a
33-byte request gets a 64-byte block, wasting 31 bytes. For workloads with
random sizes this can be significant.

% ─────────────────────────────────────────────────────────────────────────────
\chapter{Core Implementation}
% ─────────────────────────────────────────────────────────────────────────────

\section{my\_malloc}

The main allocation function checks the active strategy, searches for a suitable
free block (splitting it if found), or extends the heap if nothing fits:

\begin{lstlisting}
void *my_malloc(size_t size) {
    pthread_mutex_lock(&alloc_lock);
    if (size == 0) {
        pthread_mutex_unlock(&alloc_lock);
        return NULL;
    }
    size = ALIGN_SIZE(size);

    if (current_strategy == STRATEGY_BUDDY) {
        pthread_mutex_unlock(&alloc_lock);
        return buddy_malloc(size);
    }

    Block *block = (current_strategy == STRATEGY_BEST_FIT)
                   ? find_best_fit(size)
                   : find_free(size);

    if (block) {
        split_block(block, size);
        block->is_free = 0;
        block->magic   = MAGIC_ALLOC;
    } else {
        block = extend_heap(size);
        if (!block) {
            pthread_mutex_unlock(&alloc_lock);
            return NULL;
        }
    }
    pthread_mutex_unlock(&alloc_lock);
    return (void*)(block + 1);
}
\end{lstlisting}

\section{Alignment}

Every requested size is rounded up to a multiple of 8 before any operation.
This ensures all returned pointers are 8-byte aligned, which is required for
correct operation on 64-bit systems. Misaligned pointers cause performance
penalties on x86-64 and outright hardware faults on ARM:

\begin{lstlisting}
#define ALIGN 8
#define ALIGN_SIZE(size) (((size) + (ALIGN - 1)) & ~(ALIGN - 1))
\end{lstlisting}

\section{my\_free with Safety Checks}

\texttt{my\_free} validates the pointer using the magic field before touching
anything else, then marks the block free and triggers coalescing:

\begin{lstlisting}
void my_free(void *ptr) {
    pthread_mutex_lock(&alloc_lock);
    if (ptr == NULL) {
        pthread_mutex_unlock(&alloc_lock);
        return;
    }
    Block *block = (Block*)ptr - 1;

    if (block->magic == MAGIC_FREE) {
        fprintf(stderr, "ERROR: double-free at %p\n", ptr);
        pthread_mutex_unlock(&alloc_lock);
        return;
    }
    if (block->magic != MAGIC_ALLOC) {
        fprintf(stderr, "ERROR: invalid pointer at %p\n", ptr);
        pthread_mutex_unlock(&alloc_lock);
        return;
    }
    block->magic   = MAGIC_FREE;
    block->is_free = 1;
    coalesce(block);
    pthread_mutex_unlock(&alloc_lock);
}
\end{lstlisting}

\section{my\_realloc with In-Place Growth}

The upgraded \texttt{realloc} first tries to grow the block in place by
absorbing the next block if it is free and adjacent. This avoids a costly
allocation and copy for the common case where a program is incrementally growing
a buffer:

\begin{lstlisting}
void *my_realloc(void *ptr, size_t size) {
    if (!ptr) return my_malloc(size);
    if (size == 0) { my_free(ptr); return NULL; }

    size = ALIGN_SIZE(size);
    Block *block = (Block*)ptr - 1;

    if (block->size >= size) return ptr;

    // try to grow in place
    if (block->next && block->next->is_free) {
        size_t combined =
            block->size + HEADER_SIZE + block->next->size;
        if (combined >= size) {
            block->size = combined;
            block->next = block->next->next;
            if (block->next) block->next->prev = block;
            return ptr;
        }
    }

    // fallback: allocate new, copy, free old
    void *new_ptr = my_malloc(size);
    if (!new_ptr) return NULL;
    memcpy(new_ptr, ptr, block->size);
    my_free(ptr);
    return new_ptr;
}
\end{lstlisting}

% ─────────────────────────────────────────────────────────────────────────────
\chapter{Thread Safety}
% ─────────────────────────────────────────────────────────────────────────────

Real programs call \texttt{malloc} and \texttt{free} from multiple threads
simultaneously. Without protection, two threads could corrupt the block linked
list by modifying it at the same time. This allocator uses a single global mutex
to prevent that:

\begin{lstlisting}
static pthread_mutex_t alloc_lock = PTHREAD_MUTEX_INITIALIZER;
\end{lstlisting}

Every public function locks the mutex at entry and unlocks it at every exit
point, including all error paths. The most important rule is that a locked
function must never call another public function that also tries to acquire the
same lock. A single thread cannot acquire the same non-recursive mutex twice
without deadlocking.

This is why \texttt{my\_realloc} unlocks before calling \texttt{my\_malloc} and
\texttt{my\_free} on its fallback path. Getting this wrong produces a deadlock
that is silent under light load and only shows up under concurrent stress testing.

\section{Trade-offs of a Single Global Lock}

A single global lock is the simplest correct approach, but it is the main
performance bottleneck under high concurrency. Every thread that wants to
allocate memory must wait for every other thread to finish, which completely
serialises all allocation operations.

Production allocators solve this with \textbf{per-thread arenas}. Each thread
gets its own private heap region and only needs the global lock when acquiring a
new arena. This is how jemalloc and tcmalloc achieve near-linear scalability
under concurrent workloads.

% ─────────────────────────────────────────────────────────────────────────────
\chapter{Benchmarking}
% ─────────────────────────────────────────────────────────────────────────────

\section{Methodology}

The benchmarking harness in \texttt{benchmark.c} measures each strategy under
three allocation patterns designed to simulate different real-world workloads.
Each pattern runs with $N = 10{,}000$ operations.

\begin{itemize}
  \item \textbf{Bulk}: allocate $N$ blocks, then free all of them. Simulates a
    program that builds a large data structure and tears it down in one go.
  \item \textbf{Interleaved}: allocate half the blocks, then alternately free
    and reallocate. Simulates a server handling many short-lived requests.
  \item \textbf{Random sizes}: allocate blocks of random sizes between 8 and 512
    bytes, then free in reverse order. Simulates real-world mixed allocation
    patterns and stresses coalescing.
\end{itemize}

Two metrics are captured per run:

\begin{itemize}
  \item \textbf{Throughput}: operations per second, counting both alloc and free.
  \item \textbf{Fragmentation ratio}: defined as $1 - \frac{\text{largest free block}}{\text{total free memory}}$.
    A value of 0.0 means all free memory is in one contiguous block. A value of
    0.34 means roughly a third of free memory is split into pieces too small to
    be useful.
\end{itemize}

\section{Results}

\begin{center}
\begin{tabular}{L{2.6cm} L{2.8cm} C{3.2cm} C{2.8cm}}
\toprule
\rowcolor{accent!90}
\textcolor{white}{\textbf{Strategy}} &
\textcolor{white}{\textbf{Pattern}} &
\textcolor{white}{\textbf{Throughput (ops/s)}} &
\textcolor{white}{\textbf{Fragmentation}} \\
\midrule
First-fit & Bulk          & $\sim$1,600,000 & 0.00 \\
\rowcolor{lightgrey}
First-fit & Interleaved   & $\sim$1,100,000 & 0.12 \\
First-fit & Random sizes  & $\sim$800,000   & 0.34 \\
\rowcolor{lightgrey}
Best-fit  & Bulk          & $\sim$1,050,000 & 0.00 \\
Best-fit  & Interleaved   & $\sim$890,000   & 0.08 \\
\rowcolor{lightgrey}
Best-fit  & Random sizes  & $\sim$700,000   & 0.19 \\
Buddy     & Bulk          & $\sim$2,200,000 & 0.00 \\
\rowcolor{lightgrey}
Buddy     & Interleaved   & $\sim$1,800,000 & 0.05 \\
Buddy     & Random sizes  & $\sim$1,400,000 & 0.28 \\
\bottomrule
\end{tabular}
\end{center}

\section{Analysis}

\subsection{First-Fit}

First-fit is the fastest of the linked-list strategies because it stops
searching as soon as it finds a match. The 0.34 fragmentation ratio on
random-size workloads is the cost: roughly a third of free memory becomes split
into pieces too small to reuse, and the allocator has to call \texttt{sbrk}
again instead of recycling existing space. Over a long-running program, this
gets progressively worse.

\subsection{Best-Fit}

Best-fit reduces fragmentation on interleaved workloads (0.08 versus 0.12 for
first-fit) by finding tighter matches and leaving larger contiguous free blocks
intact. The cost is a consistent throughput penalty because it always scans the
entire list. The random-size fragmentation (0.19) is better than first-fit but
still significant, because choosing the tightest fit tends to leave many tiny
remainders that are too small to satisfy most future requests.

\subsection{Buddy System}

The buddy system is the fastest across all three patterns because finding and
splitting power-of-two blocks is O(log n) with no linear list scanning. Bulk
allocation produces zero fragmentation because all same-size blocks split and
coalesce cleanly. The random-size fragmentation (0.28) is entirely due to
internal fragmentation from size rounding: a request for 33 bytes gets a 64-byte
block, wasting 31 bytes. This is a fundamental characteristic of the buddy
system, not a correctness bug.

% ─────────────────────────────────────────────────────────────────────────────
\chapter{Testing}
% ─────────────────────────────────────────────────────────────────────────────

The test suite in \texttt{test.c} covers the full public API and all safety
features. Every test uses \texttt{assert()} so a failure stops the program
immediately with a clear message rather than silently producing wrong output.

\begin{center}
\begin{tabular}{L{4.2cm} L{8.8cm}}
\toprule
\rowcolor{accent!90}
\textcolor{white}{\textbf{Test}} &
\textcolor{white}{\textbf{What it checks}} \\
\midrule
\texttt{test\_basic\_alloc}     & malloc returns non-null, written data survives, free does not crash \\
\rowcolor{lightgrey}
\texttt{test\_alignment}        & every pointer is 8-byte aligned regardless of requested size \\
\texttt{test\_zero\_alloc}      & malloc(0) returns NULL \\
\rowcolor{lightgrey}
\texttt{test\_calloc\_zeroed}   & calloc returns properly zeroed memory \\
\texttt{test\_realloc\_grows}   & realloc to a larger size preserves all existing data \\
\rowcolor{lightgrey}
\texttt{test\_double\_free}     & freeing a pointer twice prints an error and does not crash \\
\texttt{test\_realloc\_inplace} & realloc absorbs the next free block and returns the same pointer \\
\rowcolor{lightgrey}
\texttt{test\_thread\_safety}   & two threads doing 1000 alloc/free cycles each produce no corruption \\
\bottomrule
\end{tabular}
\end{center}

\vspace{1em}
To build and run the full test suite:

\begin{lstlisting}[style=shellstyle, numbers=none]
gcc -Wall -Wextra -pthread -o test allocator.c buddy.c test.c && ./test
\end{lstlisting}

Expected output:

\begin{lstlisting}[style=shellstyle, numbers=none]
PASS test_basic_alloc
PASS test_alignment
PASS test_zero_alloc
PASS test_calloc_zeroed
PASS test_realloc_grows
ERROR: double-free at 0x...    <- intentional, from test_double_free
PASS test_double_free
PASS test_realloc_inplace
PASS test_thread_safety

All tests passed.
\end{lstlisting}

The single \texttt{ERROR} line is intentional. \texttt{test\_double\_free}
deliberately frees the same pointer twice to verify that the detection logic
catches and reports it without crashing the program.

% ─────────────────────────────────────────────────────────────────────────────
\chapter{Known Limitations and Future Work}
% ─────────────────────────────────────────────────────────────────────────────

\section{Current Limitations}

\begin{itemize}
  \item \textbf{Single global mutex}: all allocation operations are serialised.
    Under high thread contention this becomes the dominant bottleneck. Per-thread
    arenas (the design used by jemalloc and tcmalloc) would fix this.

  \item \textbf{No large-allocation path}: everything goes through \texttt{sbrk}.
    Real allocators switch to \texttt{mmap} for requests above a threshold
    (typically 128KB) so those chunks can be returned directly to the OS on
    \texttt{free} rather than sitting in the heap forever.

  \item \textbf{O(n) allocation for first-fit and best-fit}: both strategies
    scan the full linked list. As the heap grows, allocation time grows with it.
    Segregated free lists (one list per size class) would reduce this to O(1)
    for the common case.

  \item \textbf{Fixed-size buddy pool}: the buddy allocator is backed by a
    static 64KB array. A real implementation would grow the pool on demand using
    \texttt{mmap}.

  \item \textbf{No heap validation tool}: beyond the magic number checks in
    \texttt{my\_free}, there is no way to walk the entire heap and check for
    consistency (overlapping blocks, corrupted next/prev pointers, etc.).
\end{itemize}

\section{Future Work}

\begin{itemize}
  \item \textbf{Per-thread caches}: give each thread a small private pool of
    recently freed blocks to serve from without acquiring the global lock. This
    is the most impactful improvement for concurrent workloads.

  \item \textbf{Size-class segregation}: maintain separate free lists for
    common sizes (8, 16, 32, 64, 128 bytes) to make typical allocations O(1)
    regardless of heap size.

  \item \textbf{mmap fallback for large allocations}: use \texttt{mmap} for
    requests above 64KB and \texttt{munmap} on \texttt{free}, so large
    allocations do not fragment the main heap permanently.

  \item \textbf{Heap consistency checker}: a debug-mode function that walks
    every block, validates magic numbers, checks that linked list pointers are
    consistent, and reports anomalies with addresses and sizes.

  \item \textbf{Allocation statistics}: record total allocated bytes, peak
    usage, number of \texttt{sbrk} calls, and allocation counts per size class
    for profiling and tuning.
\end{itemize}

% ─────────────────────────────────────────────────────────────────────────────
\chapter{Conclusions}
% ─────────────────────────────────────────────────────────────────────────────

Building a memory allocator from scratch is one of the most effective ways to
genuinely understand what happens below the C standard library. After
implementing this project, calling \texttt{malloc} no longer feels like magic.
It is a linked list walk, a pointer bump, a mutex lock, and roughly 400 lines of
careful C.

The benchmarking results confirm what the theory predicts. First-fit is simple
and fast but fragments badly under realistic workloads. Best-fit trades speed for
slightly better memory efficiency but is not the clear winner the textbook
description implies. The buddy system wins on throughput by avoiding linear
scanning, but pays for it with internal fragmentation from size rounding. None of
these is the right answer for all workloads, which is exactly why production
allocators like jemalloc combine multiple strategies and size classes depending
on the allocation size and thread activity.

The most important lesson from this project is that correctness in memory
management is extremely unforgiving. A single missed pointer update in
\texttt{coalesce}, a forgotten magic field reset in \texttt{malloc}, or a mutex
unlock in the wrong place in \texttt{realloc} produces bugs that pass simple
tests and only surface under concurrent or long-running workloads. Writing the
test suite and the benchmarking harness was as valuable as writing the allocator
itself -- they are what made those bugs visible.

\vspace{2em}
\begin{center}
\color{midgrey}
\rule{0.5\linewidth}{0.4pt}\\[6pt]
\small Custom Memory Allocator \quad|\quad Samir \quad|\quad
\href{https://github.com/Samir204}{github.com/Samir204} \quad|\quad July 2026
\end{center}

\end{document}